\scanpage{71}
% archive foliation: 70

\begin{center}- 25 -\end{center}

\[
\left.
\begin{aligned}
H_i(0;\mu_1,\ldots,\mu_q)&=0\\
mH_i(m;\mu_1,\ldots,\mu_q)&=
\sum_{\unsure{n,(r_j),(\lambda_\ell)_\ell}}
G_i(n;r_1,\ldots,r_p;\lambda_1,\ldots,\lambda_q)\cdot\\
&\quad\cdot\prod_{j=1}^{p}\prod_{k=1}^{r_j}
H_{\unsure{j|i}}(s_{kj};\nu_{kj1},\ldots,\nu_{kjq})
\end{aligned}
\right\}\quad
\begin{gathered}(17)\\\text{pour }m>0.\end{gathered}
\]
la sommation étant étendue aux indices satisfaisant aux mêmes conditions (16) que dans (15).
La comparaison de (15) et (17) montre par récurrence sur le degré total que l'on a
\[
|h_i(m;\mu_1,\ldots,\mu_q)|\leq H_i(m;\mu_1,\ldots,\mu_q)
\]
pour tout système d'indices ; en d'autres termes, les séries
$(H_i(m;\mu_1,\ldots,\mu_q)t^m w_1^{\mu_1}\ldots w_q^{\mu_q})$
sont des \uline{majorantes} pour les séries
$(h_i(m;\mu_1,\ldots,\mu_q)x^m z_1^{\mu_1}\ldots z_q^{\mu_q})$
respectivement $(1\leq i\leq p)$.
Tout revient encore à prouver qu'on peut choisir les $G_i(n;r_1,\ldots,r_p;\lambda_1,\ldots,\lambda_q)$
de sorte que les séries majorantes
$(H_i(m;\mu_1,\ldots,\mu_q)t^m w_1^{\mu_1}\ldots w_q^{\mu_q})$
soient sommables dans un voisinage de $0$.

Or, la série entière dont la somme est fonction des $p+q+1$ variables réelles $t$, $v_i$ $(1\leq i\leq p)$ et $w_\ell$ $(1\leq\ell\leq q)$
\[
\frac{M}{(1-\alpha t)(1-\alpha\sum_{i=1}^{p}v_i)(1-\alpha\sum_{\ell=1}^{q}w_\ell)}
\]
est une majorante de chacune des séries (14) dès que $M$ est assez grand et $\alpha$ assez petit.
Tout revient à prouver qu'il existe $p$ séries entières sans terme indépendant de $t$, qui satisfont aux équations
\[
(18)\qquad
\frac{dv_k}{dt}=
\frac{M}{(1-\alpha t)(1-\alpha\sum_{i=1}^{p}v_i)(1-\alpha\sum_{\ell=1}^{q}w_\ell)}
\qquad(1\leq k\leq p)
\]
On déduit aussitôt de (18) que si ce système admet une solution, tous les $v_k$ sont égaux à $v_1=v$ satisfaisant à l'équation différentielle
